\title{xLSTM: Extended Long Short-Term Memory}

\begin{abstract}
The foundational concepts behind Long Short-Term Memory (LSTM)—namely, the constant error carousel and gating mechanisms—originated in the 1990s. LSTMs have since persisted as a staple in deep learning, playing a pivotal role in the emergence of the first Large Language Models (LLMs). Despite these achievements, the Transformer architecture, powered by efficient, parallelizable self-attention, has recently outstripped LSTMs in large-scale applications, initiating a transformative shift in the field. This work asks a straightforward question: What are the attainable results in language modeling if we scale LSTMs to billions of parameters, incorporating contemporary advancements from LLMs, while addressing the longstanding shortcomings of LSTMs? To this end, we present exponential gating enhanced by appropriate normalization and stabilization strategies. We also redesign LSTM’s memory components, resulting in: (i) sLSTM, featuring a scalar memory with scalar updates and an updated memory integration process, and (ii) mLSTM, characterized by a matrix-based memory and a covariance-inspired update rule that allows for full parallelization. By embedding these novel LSTM variants within residual blocks, we construct xLSTM blocks, which are then stacked in a residual manner to create xLSTM architectures. The combination of exponential gating and the overhauled memory mechanisms substantially enhances the expressiveness of xLSTM, enabling these architectures to rival and often surpass leading Transformers and State Space Models, both in terms of scaling properties and predictive performance.\\
Code available at: \url{https://github.com/NX-AI/xlstm}
\end{abstract}

\section{Introduction}
\label{sec:intro}

The Long Short-Term Memory (LSTM) ideas~\citep{Hochreiter:91,Hochreiter:97nips,Hochreiter:97}, i.e., the constant error carousel and gating, were introduced to overcome the vanishing gradient problem of recurrent neural networks~\citep{Hochreiter:91,Hochreiter:00book}:

\begin{align}
\label{eq:lstm_idea}
  c_t \ = \ f_t \ c_{t-1} \ + \ 
          i_t \  z_t \ , \quad  
          h_t \ = \ o_t \ \psi({c_t}) \ . 
\end{align}

The constant error carousel involves an additive modification of the previous cell state $c_{t-1}$ (represented in green), with contributions from cell inputs $z_t$ and regulation through sigmoid gating mechanisms (shown in blue). This updating process is governed by both the input gate $i_t$ and the forget gate $f_t$, whereas the output gate $o_t$ determines the release of the memory cell’s output, specifically the hidden state $h_t$. After being transformed or normalized via $\psi$, the cell state undergoes output gating, resulting in the final hidden state.

LSTMs have achieved notable success across a broad spectrum of application areas \citep{Hochreiter:01,Hochreiter:07,Schmidhuber:15}, maintaining dominance in text generation prior to the emergence of Transformers in 2017~\citep{Vaswani:17}. The versatility and power of LSTMs have been empirically established through their application to diverse sequence-based problems, including, but not limited to, text generation~\citep{Graves:13,Karpathy:15blog}, handwritten text synthesis~\citep{Graves:13}, sequence-to-sequence translation tasks~\citep{Sutskever:14nips}, analysis of computer programs~\citep{Zaremba:14arxiv}, the production of descriptive image captions~\citep{Karpathy:15,Hossain:19}, source code generation~\citep{Karpathy:15blog}, hydrological tasks such as rainfall-runoff prediction~\citep{Kratzert:18,Kratzert:19}, and operational flood warning models~\citep{Nearing:24}. Within the realm of reinforcement learning, LSTMs currently represent the leading architecture for sequential data, as exemplified by their adoption in systems such as the AlphaStar agent for StarCraft II~\citep{Vinyals:17short}, OpenAI Five for Dota 2~\citep{OpenAI:19}, and in nuclear fusion magnetic controller models~\citep{Degrave:22short}. Notably, LSTMs demonstrate considerable proficiency in learning and representing abstractions—effectively capturing semantic features within their memory components~\citep{Karpathy:15blog}—a property evidenced by discoveries of specific neurons attuned to numerical, syntactic~\citep{Lakretz:19}, linguistic~\citep{Bau:19}, or sentiment~\citep{Radford:17arxiv} cues. Even today, LSTMs maintain their relevance and robustness across critical applications~\citep{Degrave:22short,Nearing:24}, attesting to their enduring value.

Although LSTMs have achieved significant advancements, they possess three key shortcomings: (i) Inability to alter previous storage choices. We illustrate this through the {\it Nearest Neighbor Search} task (refer also to Appendix~\ref{sec:appNearestNeighborSearch}): When a query vector is provided, the sequence is processed one step at a time to locate the most similar vector and retrieve its associated value at the sequence's conclusion. The mean squared error for this task appears in the left panel of Figure~\ref{fig:lstmProblems}. LSTMs exhibit difficulty in updating a stored entry if a superior match is subsequently encountered. In contrast, the newly introduced xLSTM addresses this by employing exponential gating. (ii) Constrained memory capacity, as information is reduced to scalar-valued cell states. This restriction is highlighted with {\it Rare Token Prediction}. In the right panel of Figure~\ref{fig:lstmProblems}, we present the perplexity on Wikitext-103~\citep{Merity:17} for groups based on token rarity. The LSTM's performance is hindered for infrequent tokens due to its limited capacity, an issue xLSTM overcomes with a matrix-based memory. (iii) Sequential computation inherent in LSTMs, caused by memory mixing, namely the recurrence between subsequent hidden states that prevents parallel computation.

These deficiencies in LSTM architectures have contributed to the widespread adoption of Transformers~\citep{Vaswani:17} for language modeling tasks. An open question arises: What levels of performance in language modeling are attainable when these LSTM constraints are eliminated and the architecture is scaled comparably to modern Large Language Models?

\section{Extended Long Short-Term Memory}
\label{sec:xLSTM}

To address the limitations inherent in traditional LSTM, Extended Long Short-Term Memory (xLSTM) incorporates two primary enhancements to the framework outlined in Equation~\eqref{eq:lstm_idea}. These advancements---exponential gating mechanisms and the introduction of novel memory architectures---add two distinct members to the LSTM family: (i) sLSTM (detailed in Section~\ref{sec:sLSTM}), which operates with scalar memory, leverages scalar updates, and involves memory mixing; and (ii) mLSTM (explained in Section~\ref{sec:mLSTM}), which uses a matrix-based memory along with a covariance update rule, based on the outer product, supporting complete parallelization. Both of these, sLSTM and mLSTM, benefit from the exponential gating innovation. To further facilitate parallel computation, mLSTM omits memory mixing (that is, it does not include hidden-to-hidden recurrent connections). In both architectures, multiple memory cells are feasible: sLSTM allows mixing between memory cells, whereas mLSTM treats memory cells independently. Additionally, sLSTM accommodates multiple heads, where each head does not engage in mixing with other heads but maintains mixing among cells within a single head. This combination of multiple heads and exponential gating introduces a novel paradigm for memory mixing in sLSTM. In contrast, for mLSTM, the concepts of having multiple heads and multiple cells coincide.

Embedding these new sLSTM and mLSTM variants as modular elements within residual blocks generates the xLSTM blocks (described in Section~\ref{sec:xLSTMblock}). A series of these blocks, when stacked residually within a network, form an xLSTM-based architecture (see Section~\ref{sec:xLSTMarchitecure} for details). The architectural design and its constituent parts are illustrated in Figure~\ref{fig:xlstm_sketch}.

\subsection{Review of the Long Short-Term Memory}
\label{sec:orgLSTM}

The foundational concept behind LSTM~\citep{Hochreiter:91,Hochreiter:97nips,Hochreiter:97} centers on the scalar memory cell, which functions as both a storage and processing mechanism. This architecture ensures persistent gradients, circumventing the issue of vanishing gradients~\citep{Hochreiter:91,Hochreiter:00book}, primarily due to the constant error carousel inherent in the cell state update procedure. Within the memory cell, three distinct gates are present: the input gate, the output gate, and the forget gate—this last element was incorporated later by \citet{Gers:00}. The LSTM memory cell's update equations at time step $t$ are outlined as follows:

\begin{align}
c_t \ &= \  f_t \ c_{t-1} \ + \ i_t \ z_t &  & &\text{cell state} \\
h_t  \ &= \ o_t \ \tilde{h}_t \ , & \tilde{h}_t \ &= \ \psi \left( c_t\right)
  &\text{hidden state} \\
z_t \ &= \ \varphi \left( \tilde{z}_t \right) \ , 
  &\tilde{z}_t \ &=  \ \mathbf{w}^\top_{z} \ \mathbf{x}_t \ + \
  r_{z}  h_{t-1} \ + \  b_{z} \ \
  &\text{cell input} \\
i_t \ &= \ \sigma \left( \tilde{i}_t  \right) \ , 
  &\tilde{i}_t \ &= \ \mathbf{w}^\top_{i} \ \mathbf{x}_t \ + \
  r_{i}  \ h_{t-1} \ + \  b_{i} \ \
  &\text{input gate} \\
f_t \ &= \ \sigma \left(  \tilde{f}_t \right) \ , 
  &\tilde{f}_t \ &= \ \mathbf{w}^\top_{f} \ \mathbf{x}_t  \ + \
  r_{f}  \ h_{t-1} \ + \  b_{f} \ \
  &\text{forget gate} \\
o_t \ &= \ \sigma \left( \tilde{o}_t \right) \ , 
  &\tilde{o}_t  \ &= \ \mathbf{w}^\top_{o} \ \mathbf{x}_t \ + \
  r_{o}  \ h_{t-1} \ + \  b_{o} \ \
  &\text{output gate} 
\end{align}

The vectors $\mathbf{w}_{z}$, $\mathbf{w}_{i}$, $\mathbf{w}_{f}$, and $\mathbf{w}_{o}$ represent the input weight parameters connecting the input vector $\mathbf{x}_t$ to the cell input, as well as the input, forget, and output gates, respectively. The parameters $r_{z}$, $r_{i}$, $r_{f}$, and $r_{o}$ denote the recurrent weight connections from the previous hidden state $h_{t-1}$ to the corresponding components: the cell input, input gate, forget gate, and output gate. Corresponding biases are denoted as $b_{z}$, $b_{i}$, $b_{f}$, and $b_{o}$. The functions $\varphi$ and $\psi$ are used as activation functions for the cell input and hidden state, commonly instantiated as $\tanh$. The activation function $\psi$ serves to constrain the range of the cell state, which otherwise could grow without bounds. The gates utilize sigmoid activation functions defined as $\sigma \left( x \right)= 1/(1+\exp(-x))$. Subsequent versions unified multiple scalar-valued memory cells as a vector $\mathbf{c}_t \in \mathbb{R}^{d}$, enabling the introduction of recurrent weight matrices $\mathbf{R} \in \mathbb{R}^{d \times d}$ that allow interactions among outputs of individual memory cells~\citep{Greff:15}; further information can be found in Appendix~\ref{sec:apporgLSTM}. Studies involving the removal of model components established that every part of the memory cell plays an essential role~\citep{Greff:15}.

\subsection{sLSTM}
\label{sec:sLSTM}

We enhance LSTMs by enabling them to modify storage choices through the incorporation of exponential gates (red), coupled with normalization and stabilization mechanisms. Specifically, exponential activation functions are employed for both input and forget gates. For the normalization component, a normalizer state is defined, which accumulates the products of the input gate and all subsequent forget gates.

The scalar sLSTM forward pass is:

\begin{align}
\label{eq:slstmforward}
c_t \ &= \  f_t \ c_{t-1} \ + \ i_t \ z_t & & 
 &\text{cell state} \\
n_t \ &= \  f_t \ n_{t-1} \ + \ i_t  & & 
 &\text{normalizer state} \\
h_t \ &= \ o_t \ \tilde{h}_t \ ,
  &\tilde{h}_t \ &= \ c_t / n_t
&\text{hidden state} \\
z_t \ &= \ \varphi \left( \tilde{z}_t \right) \ , 
  &\tilde{z}_t \ &=  \ \mathbf{w}^\top_{z} \ \mathbf{x}_t \ + \
  r_{z}  h_{t-1} \ + \  b_{z}
  &\text{cell input} \\
i_t \ &= \ \!\exp \left( \tilde{i}_t  \right) \ , 
  &\tilde{i}_t \ &= \ \mathbf{w}^\top_{i} \ \mathbf{x}_t \ + \
  r_{i}  \ h_{t-1} \ + \  b_{i}
  &\text{input gate} \\
f_t \ &= \ \sigma \left(  \tilde{f}_t \right) \ \text{OR} \
\!\exp \left(  \tilde{f}_t \right) \ ,
  &\tilde{f}_t \ &= \ \mathbf{w}^\top_{f} \ \mathbf{x}_t  \ + \
  r_{f}  \ h_{t-1} \ + \  b_{f}
  &\text{forget gate} \\
o_t \ &= \ \sigma \left( \tilde{o}_t \right) \ , 
  &\tilde{o}_t  \ &= \ \mathbf{w}^\top_{o} \ \mathbf{x}_t \ + \
  r_{o}  \ h_{t-1} \ + \  b_{o}
  &\text{output gate} 
\end{align}

We transfer the original LSTM gating techniques, i.e., input- and/or hidden-dependent gating plus bias term, to the new architectures. Exponential activation functions can lead to large values that cause overflows. Therefore, we stabilize gates with an additional state \mbox{$m_t$~\citep{Milakov:18arxiv}}:

\begin{align}
\label{eq:slstmstabil}
m_t \ &= \ \max \left( \log ( f_t ) + m_{t-1} , \log ( i_t ) \right) &\text{stabilizer state} \\
i'_t \ &= \ \!\exp \left( \log \left ( i_t \right) - m_t \right) \ = \ \!\exp \left( \tilde{i}_t - m_t \right) \ \, 
  &\text{stabil. input gate} \\
f'_t \ &= \ \!\exp \left( \log \left( f_t \right) + m_{t-1} - m_t \right) \ \, 
  &\text{stabil. forget gate}
\end{align}

Appendix~\ref{sec:appsLSTM} demonstrates that substituting $f_t$ with $f'_t$ and $i_t$ with $i'_t$ during the forward pass leaves both the network’s output and the loss function’s parameter derivatives unaffected.

\paragraph{New Memory Mixing.} Similarly to the standard LSTM, sLSTM architecture may employ several memory cells, as detailed in Appendix~\ref{sec:appsLSTM}. Utilizing multiple memory cells facilitates the mixing of information through the recurrence matrices $\mathbf{R}_{\mathbf{z}}$, $\mathbf{R}_{\mathbf{i}}$, $\mathbf{R}_{\mathbf{f}}$, $\mathbf{R}_{\mathbf{o}}$, which transmit signals from the hidden state vector $\mathbf{h}$ to the cell input $\mathbf{z}$ and to each gating mechanism—$\mathbf{i}$, $\mathbf{f}$, and $\mathbf{o}$—accordingly. Exponential gating adds an additional dimension to this memory mixing process. The updated sLSTM model supports several heads, enabling memory mixing within individual heads, yet not between different heads. By combining the notion of heads and exponential gating in sLSTM, a distinct paradigm for memory mixing emerges.

\subsection{mLSTM}
\label{sec:mLSTM}

To improve the memory capacity of LSTMs, we replace the scalar memory cell $c \in \mathbb{R}$ with a matrix cell $\mathbf{C} \in \mathbb{R}^{d \times d}$. This modification allows retrieval operations to use matrix multiplication. At a particular timestep $t$, our aim is to store a vector pair consisting of a key $\mathbf{k}_t \in \mathbb{R}^d$ and a value $\mathbf{v}_t \in \mathbb{R}^d$, following Transformer nomenclature. Subsequently, at timestep $t + \tau$, we use a query vector $\mathbf{q}_{t+\tau} \in \mathbb{R}^d$ to recover the associated value $\mathbf{v}_t$. This corresponds to the framework utilized in Bidirectional Associative Memories (BAMs) \citep{Kohonen:72,Anderson:72,Nakano:72,Anderson:77}.

The covariance update rule~\citep{Sejnowski:77,Dayan:91} for storing a key-value pair is

\begin{align}
\mathbf{C}_t \ &= \ \mathbf{C}_{t-1} \ + \ \mathbf{v}_{t} \ \mathbf{k}_t^\top \ .
\end{align}

We presume the presence of layer normalization preceding the projection of inputs to keys and values, ensuring these have a mean of zero. As established in prior work~\citep{Dayan:91}, updating the covariance matrix in this manner achieves optimal conditions for maximizing the separability among retrieved binary vectors; this optimization is tantamount to maximizing the signal-to-noise ratio. Restricting retrieval to pairwise associations while accepting quadratic computational demands, as seen in attention mechanisms~\citep{Krotov:16,Krotov:17,Ramsauer:21}, can further enhance separability. The covariance update mechanism also aligns with the approach used in Fast Weight Programmers~\citep{Schmidhuber:92ncfastweights,Schlag:21}, where more recent implementations incorporate a decay factor applied to $\mathbf{C}_{t-1}$ and a learning factor for updating with 
$\mathbf{v}_{t} \mathbf{k}_t ^\top$~\citep{Ba:16}. Taking inspiration from this, we embed the covariance update process into the LSTM architecture, mapping the forget gate to the decay factor, the input gate to the learning rate, and utilizing the output gate to rescale the retrieved output.

For this matrix memory, the normalizer state is the weighted sum of key vectors, where each key vector is weighted by the input gate and all future forget gates. Again, the normalizer state keeps record of the strength of the gates. Since the dot product between query and normalizer state can be close to zero, we use the absolute value of this dot product and lower bound it by a threshold (typically 1.0) as done previously~\citep{Sun:23arxiv}. The mLSTM forward pass is:

\begin{align}
\label{eq:mlstm_recurrent_begin}
\mathbf{C}_t \ &= \  f_t \ \mathbf{C}_{t-1} \ + \ 
  i_t \ \mathbf{v}_t \ \mathbf{k}_t^\top & &  &\text{cell state} \\
\mathbf{n}_t \ &= \  f_t \ \mathbf{n}_{t-1} \ + \ 
  i_t \ \mathbf{k}_t & &  &\text{normalizer state} \\
\mathbf{h}_t  \ &= \ \mathbf{o}_t \ \odot \ \tilde{\mathbf{h}}_t
\ , \qquad \qquad 
\tilde{\mathbf{h}}_t \ = \   \mathbf{C}_t \mathbf{q}_t \ / \ 
  \max \left\{ |\mathbf{n}_t^\top \mathbf{q}_t|, 1 \right\} 
   &&&\text{hidden state} \\
\mathbf{q}_t \ &= \ \mathbf{W}_q \ \mathbf{x}_t \ + \ \mathbf{b}_q  & &  &\text{query input} \\
\mathbf{k}_t \ &= \ \frac{1}{\sqrt{d}} \mathbf{W}_k \ \mathbf{x}_t \ + \ \mathbf{b}_k  & &  &\text{key input} \\
\mathbf{v}_t \ &= \ \mathbf{W}_v \ \mathbf{x}_t \ + \ \mathbf{b}_v  & &  &\text{value input} \\
i_t \ &= \ \!\exp \!  \! \left( \tilde{i}_t  \right) \ , 
 \qquad \qquad \ \ \ \ \,
  \tilde{i}_t \ = \ \mathbf{w}^\top_{i} \ \mathbf{x}_t \ + \  b_{i} \ \ \ 
  &&&\text{input gate} \\
f_t \ &= \ \sigma \!  \left(  \tilde{f}_t \right) \ \text{OR} \ \!\exp \!  \! \left(  \tilde{f}_t \right) , \ \ \: \!
\tilde{f}_t \ = \ \mathbf{w}^\top_{f} \ \mathbf{x}_t  \ + \
  b_{f}
  &&& \text{forget gate} \\
\label{eq:mlstm_recurrent_end}
\mathbf{o}_t \ &= \ \sigma \left( \tilde{\mathbf{o}}_t \right) \ , \qquad \qquad \quad \ \ \ \,
  \tilde{\mathbf{o}}_t  \ = \ \mathbf{W}_{\mathbf{o}} \ \mathbf{x}_t \ + \
  \mathbf{b}_{\mathbf{o}}
  &&&\text{output gate} 
\end{align}

Similar to standard LSTM architectures, mLSTM accommodates several memory cells. In the case of mLSTM, having multiple heads is functionally the same as having multiple cells because memory states remain distinct and do not interact. To mitigate the instability that arises from the exponential gates within mLSTM, the stabilization strategies established for sLSTM are applied here as well (refer to Equation~\ref{eq:slstmstabil}). Due to the absence of memory state intermixing in mLSTM, the recurrence relation is amenable to restructuring for parallel computation. Additional information is provided in Appendix~\ref{sec:appmLSTM}.

\subsection{xLSTM Architecture}
\label{sec:xLSTMarch}

\paragraph{xLSTM Blocks.}
\label{sec:xLSTMblock}
The purpose of an xLSTM block is to form a nonlinear representation of the preceding sequence within a higher-dimensional space, facilitating the distinction of diverse histories or contextual backgrounds. This ability to distinguish between pasts is essential for accurately forecasting the subsequent element in a sequence, such as predicting the following token. Our approach leverages Cover’s Theorem~\citep{Cover:65}, which suggests that patterns that are nonlinearly embedded into a space of greater dimensionality become more likely to be linearly separable in comparison to the original space. We examine two types of residual block structures: (i) a residual block employing a post up-projection mechanism, similar to what is used in Transformers. In this case, the past is first nonlinearly encoded within the original dimension, subsequently linearly projected into a higher-dimensional space, processed through a nonlinear activation, and finally projected back to the initial space via a linear operation; this is illustrated in the left panel of Figure~\ref{fig:Backbones} and in the third column of Figure~\ref{fig:xlstm_sketch}. A more comprehensive illustration is available in Figure~\ref{fig:appsLSTM_detailed} located in the appendix. (ii) a residual block that uses a pre up-projection strategy, akin to what is found in State Space Models, which first applies a linear map into a high-dimensional space,

Non-linear summarization of previous states occurs in the high-dimensional representation space, which is then linearly projected back into the original dimensionality. In xLSTM blocks incorporating sLSTM, the post up-projection block is employed. Conversely, for xLSTM blocks with mLSTM, we adopt the pre up-projection block, as the higher-dimensional space accommodates increased memory capacity. For further specifics, consult the leftmost panel of Figure~\ref{fig:Backbones}, the third column in Figure~\ref{fig:xlstm_sketch}, or Figure~\ref{fig:appsLSTM_detailed} in the appendix.

\paragraph{xLSTM Architecture. }
\label{sec:xLSTMarchitecure}
The xLSTM architecture is realized by stacking its component blocks with residual connections, as described in~\citet{Srivastava:15,He:16}. Our implementation follows the widely adopted pre-LayerNorm~\citep{Ba:16b} residual configuration, prevalent in modern Large Language Models. Refer to the final column of Figure~\ref{fig:xlstm_sketch} for an illustration.

\subsection{Memory and Speed Considerations}

Unlike Transformers, xLSTM networks exhibit linear computational demands and constant memory usage as the sequence length grows. Owing to the compressive nature of xLSTM memory, these architectures are particularly advantageous for industrial scenarios and edge deployment.

In the case of mLSTM, its memory component is parameter-free but incurs high computational cost because of the $d \times d$ matrix memory and corresponding $d \times d$ update operations. This design represents a compromise between increasing memory capacity and computational expense. However, the computations are amenable to GPU-based parallelization, which ensures only a marginal impact on actual execution time.

While mLSTM can be parallelized in a manner similar to FlashAttention~\citep{Dao:22,Dao:23} or GLA~\citep{Yang:23arxiv}, sLSTM's memory mixing—resulting from hidden-to-hidden interactions—precludes parallelism. Nonetheless, we have engineered a high-performance CUDA implementation that incorporates optimizations down to GPU register usage, typically yielding speeds less than double those of mLSTM.

\section{Related Work}
\label{sec:related}

\paragraph{Linear Attention.} Numerous strategies have been introduced to address the Transformer’s quadratic computational scaling with respect to context length, aiming to achieve linear complexity in attention. The Synthesizer framework constructs attention weights synthetically, eliminating direct token-to-token dependencies~\citep{Tay:20arxiv}. Linformer models self-attention using low-rank matrices and further provides linear approximations~\citep{Wang:20arxiv}. Linear Transformer transforms the attention operation into a linear formulation~\citep{Katharopoulos:20}. Performer utilizes a method based on positive orthogonal random features to linearly approximate the softmax used in attention~\citep{Choromanski:21}. Alternatives such as Structured Global Convolution (SGConv)~\citep{Li:22} and Hyena Hierarchy~\citep{Poli:23} have replaced traditional attention with efficient long convolution operations.

\paragraph{State Space Models.} In recent years, State Space Models (SSMs) have gained substantial traction owing to their linear complexity with respect to sequence length and their competitive results compared to Transformer architectures. The early contributions in this area include the Structured State Space sequence model (S4)~\citep{Gu:21}. This was succeeded by subsequent advancements such as the Diagonal State Space (DSS) model~\citep{Gupta:22}, Gated State Space (GSS) models~\citep{Mehta:22}, the S5 model~\citep{Smith:22}, Bidirectional Gated SSM (BiGS)~\citep{Wang:22}, H3 model~\citep{Fu:23}, as well as Mamba~\citep{Gu:24arxiv}.

\paragraph{Recurrent Neural Networks.} In recent years, Recurrent Neural Networks (RNNs) have been explored as substitutes for Transformers and attention mechanisms, primarily due to their linear scalability with respect to sequence length. Models such as RNNs with Deep Linear Recurrent Units (LRUs) have demonstrated significant advancements in language modeling~\citep{Orvieto:23, De:24arxiv}. Similarly, architectures like Hierarchically Gated Linear RNN (HGRN)~\citep{Qin:23} and its successor HGRN2~\citep{Qin:24arxiv} have shown strong results. Among notable RNN-based frameworks, RWKV~\citep{Peng:23arxivshort,Peng:24arxivshort} has emerged as a competitive alternative to Transformer models for large language modeling tasks.

\paragraph{Gating.}
Gating mechanisms, initially introduced in LSTM architectures, have since been revisited and adapted in numerous modern methods. Variants of gating appear across HGRN~\citep{Qin:23}, HGRN2~\citep{Qin:24arxiv}, Gated Linear Attention (GLA)~\citep{Yang:23arxiv}, Gated State Space (GSS) models~\citep{Mehta:22}, Bidirectional Gated SSM (BiGS)~\citep{Wang:22}, Moving Average Equipped Gated Attention (MEGA)~\citep{Ma:22}, RWKV~\citep{Peng:23arxivshort}, and Mamba~\citep{Gu:24arxiv}.

\paragraph{Covariance Update Rule.} To improve memory retention, the mLSTM cell was augmented with a matrix memory governed by a covariance update rule. Several related strategies utilize similar update schemes, including Fast Weight Programmers \citep{Schmidhuber:92ncfastweights,Schlag:21}, RWKV-5 and RWKV-6~\citep{Peng:24arxivshort}, Retention~\citep{Sun:23arxiv}, Linear Transformer~\citep{Katharopoulos:20}, and HGRN2~\citep{Qin:24arxiv}.

\paragraph{Most Related.} The models most analogous to xLSTM are Retention~\citep{Sun:23arxiv}, RWKV~\citep{Peng:23arxivshort,Peng:24arxivshort}, and HGRN2~\citep{Qin:24arxiv}, which employ matrix-based memory and often incorporate gating mechanisms. Unlike the newly proposed sLSTM, these prior architectures lack support for memory mixing. The ability to mix memory is crucial for addressing state tracking tasks, making LSTM variants more capable than State Space Models (SSMs) or Transformers~\citep{Merrill:24,Deletang:23}. State tracking plays a pivotal role in tasks such as code evaluation and monitoring entities throughout extended texts.

\paragraph{Residually Stacking Architectures.} The xLSTM family, like most state-of-the-art deep learning systems, is engineered by means of stacking its constituent modules in a residual manner~\citep{Srivastava:15,He:16}.

Employing this structural paradigm has proven fundamental in the evolution of deep convolutional neural networks~\citep{He:16} and Transformer models~\citep{Vaswani:17}. Transformer-based approaches have propelled the development of Large Language Models (LLMs), including GPT-3~\citep{Brown:20short}, ChatGPT~\citep{Schulman:22short}, GPT-4~\citep{Achiam:23gpt4short}, Megatron-LM~\citep{Shoeybi:19}, Gopher~\citep{Rae:21short}, ERNIE 3.0 Titan~\citep{Wang:21arxivshort}, GLaM~\citep{Du:21glamshort}, Chinese M6~\citep{Lin:21}, the multilingual AlexaTM 20B~\citep{Soltan:22}, OPT~\citep{Zhang:22opt}, Chinchilla~\citep{Hoffmann:22short}, BLOOM~\citep{Scao:22arxivshort}, GLM-130B~\citep{Zeng:22arxivshort}, LaMDA~\citep{Thoppilan:22short}, PaLM~\citep{Chowdhery:22short}, Llama~\citep{Touvron:23arxiv}, and Gemini~\citep{GeminiTeam:23arxiv,Reid:24arxiv}.

\section{Experiments}
\label{sec:Exp}

This section presents an empirical analysis of xLSTM and benchmarks it against prevailing approaches, with emphasis on language modeling tasks. Section~\ref{sec:ExpSyntethic} explores the distinct capabilities of xLSTM using synthetic experiments. 
Section~\ref{sec:ExpComparison} provides a comparative study of language modeling methods by evaluating their validation set perplexity after training on 15B tokens drawn from SlimPajama~\citep{Soboleva:23}. Ablation experiments with xLSTM are also conducted using this dataset. Subsequently, the scaling properties of each technique are examined in line with methodologies employed by \citet{Kaplan:20} and \citet{Brown:20short}.
A comprehensive evaluation of language modeling is performed in Section~\ref{sec:ExpLanguage}. Here, xLSTM and the top-performing approaches identified in Section~\ref{sec:ExpComparison} are trained with 300B tokens from SlimPajama~\citep{Soboleva:23}. This section proceeds to: (1) analyze each method's ability to generalize to longer contexts; (2) assess performance based on validation perplexity and downstream tasks following \citep{Sutawika:24}; (3) measure effectiveness across 571 distinct text domains in the PALOMA language benchmark dataset~\citep{Magnusson:23arxivshort}; and (4) re-examine scaling behavior for all methods, this time using training data increased by a factor of 20.

\label{sec:xlstm_blocks_notation}
Throughout all experimental setups, we denote xLSTM[$a$:$b$] to indicate a proportion $a/b$ of mLSTM-based to sLSTM-based xLSTM blocks. For instance, xLSTM[7:1] specifies that out of every eight blocks, seven are built using mLSTM and one utilizes sLSTM. Given a typical block count of 48, this results in 42 mLSTM-based blocks and 6 sLSTM-based blocks. Additionally, in every experiment, mLSTM components are paired with pre up-projection blocks while sLSTM components incorporate post up-projection blocks.

\subsection{Synthetic Tasks and Long Range Arena}
\label{sec:ExpSyntethic}
We begin by evaluating the impact of xLSTM's exponential gating mechanism with memory mixing on formal language tasks~\citep{Deletang:23}. Subsequently, we examine xLSTM's novel matrix memory approach using the Multi-Query Associative Recall benchmark~\citep{Arora:23arxiv}. Lastly, we analyze how xLSTM performs on lengthy sequence processing tasks within the Long Range Arena~\citep{Tay:21}.

\paragraph{Evaluation of xLSTM's Exponential Gating with Memory Mixing.} We assess the capability of xLSTM's novel exponential gating mechanism, which incorporates memory mixing and is designed to address state tracking challenges~\citep{Merrill:24, Merrill:23}. To facilitate length generalization, we adapt and extend the formal language benchmarks from \citet{Deletang:23} to allow training across multiple input lengths. For comprehensive explanations of each benchmark and additional experimental outcomes, refer to Appendix~\ref{sec:appExpSynthetic-formalLang}. 

The performance of xLSTM is contrasted with a range of models, including Transformers, State Space Models, and Recurrent Neural Networks. Evaluation is performed on tokens pertinent to the specific task, and accuracy is normalized from 0 (equivalent to random chance) to 1 (achieving flawless predictions). Specifically, we test 2-block variants of several architectures: xLSTM[0:1] (representing only sLSTM), xLSTM[1:0] (representing only mLSTM), xLSTM[1:1], Llama, Mamba, RWKV, Retention, Hyena, LSTM, and LSTM inserted within Transformer blocks (LSTM (Block)). Figure~\ref{fig:formal-main} displays the experimental findings.

We observe that architectures such as Transformers or State Space Models lacking explicit memory mixing (and therefore incapable of state tracking) are unable to succeed on tasks such as the parity regular grammar. This aligns with previous evidence indicating that, in contrast to RNNs, both Transformers and State Space Models are intrinsically limited in computational power for these tasks~\citep{Merrill:24, Merrill:23, Deletang:23}.

\paragraph{Test of xLSTM's Memory Capacities on Associative Recall Tasks.} In this study, we evaluate the effectiveness of the newly introduced matrix memory in xLSTM using the Multi-Query Associative Recall task~\citep{Arora:23arxiv}. Within this framework, sequences are composed of key-value pairs, each selected randomly from a comprehensive vocabulary set, and models are required to retain this information for subsequent recall. To intensify the challenge compared to the standard task, we scale the number of key-value pairs as high as 256 and lengthen the context window to 2048 tokens, thus providing an expanded benchmark for assessing memory capabilities across different architectures. Our assessment encompasses 2-block configurations of Llama, Mamba, RWKV-5, RWKV-6, xLSTM[1:1], and xLSTM[1:0], with model performance measured by the proportion of pairs correctly retrieved. Transformers such as Llama are known to possess memory capacities exponential in the representation dimension~\citep{Ramsauer:21}, establishing a high bar for this task. Figure~\ref{fig:mqar-main} reports these comparative results. Among models outside the Transformer family, xLSTM[1:1] yields the top recall accuracy, maintaining strong performance even at smaller scales. Notably, incorporating the sLSTM block does not impair memory performance; in fact, it appears to augment memory retention, particularly evident when handling sequences containing the maximal set of 256 key-value pairs. Further analyses, including studies of extrapolation to previously unseen context lengths, are provided in Appendix~\ref{sec:appExpSynthetic-mqar}, supporting the conclusion that xLSTM's advanced memory architecture accommodates effective generalization beyond training regimes.

\paragraph{Test of xLSTM's Long Context Capabilities on Long Range Arena.} In order to evaluate xLSTM's effectiveness with lengthy sequences and extended contexts, we benchmark various approaches using the Long Range Arena~\citep{Tay:21}. Across all tested tasks, xLSTM consistently achieves high performance, indicating that its architecture is highly adept at addressing a wide range of long context challenges.

For more details, see Appendix~\ref{sec:appExpSynthetic-lra}.

\subsection{Method Comparison and Ablation Study}
\label{sec:ExpComparison}

In order to investigate our central research question—namely, the performance of our proposed LSTM variants at scale in language modeling tasks—we conduct experiments by training xLSTMs, Transformers, State Space Models, and additional architectures on a corpus of 15B tokens from SlimPajama, all within a unified auto-regressive framework. Model performance is evaluated on a shared validation dataset, and we carry out ablation analyses specifically targeting the xLSTM architectures.

\paragraph{Comparative Analysis of xLSTM and Other Architectures.} To facilitate a fair comparison, all models are trained on a dataset comprising 15B tokens from SlimPajama~\citep{Soboleva:23}. Their performance is assessed based on perplexity measured against a common validation set. The evaluated architectures include: xLSTM (the novel approach introduced here), GPT-3 (Transformer)~\citep{Brown:20short}, Llama (Transformer)~\citep{Touvron:23arxiv}, H3 (SSM)~\citep{Fu:23}, Mamba (SSM)~\citep{Gu:24arxiv}, RWKV-4 (RNN)~\citep{Peng:23arxivshort}, RWKV-5 (RNN)~\citep{Peng:24arxivshort}, RWKV-6 (RNN)~\citep{Peng:24arxivshort}, GLA (linear Transformer)~\citep{Yang:23arxiv}, HGRN (RNN)~\citep{Qin:23}, HGRN2 (RNN)~\citep{Qin:24arxiv}, RetNet (linear Transformer)~\citep{Sun:23arxiv}, Hyena (linear Transformer)~\citep{Poli:23}, along with xLSTM[1:0] and xLSTM[7:1] variants (see Section~\ref{sec:xlstm_blocks_notation}). 

Mixed precision is employed in the training regimen. However, for RWKV-5, RWKV-6, GLA, and HGRN2, the mixed precision was implemented without leveraging PyTorch's automated mixed precision functionality (refer also to Appendix Section~\ref{sec:appGenTrainProc}). 

The architectures are grouped as follows: (a) Transformers, (b) State Space Models (SSMs), (c) Recurrent Neural Networks (RNNs) together with linear Transformers. Linear Transformers utilize linear operations in place of standard Transformer attention. All compared models conform to the GPT-3 scale, with a parameter count of 350M, embedding dimension of 1024, and 24 residual blocks. Notably, GPT-3 alone employs tied weights for token and output embeddings, which reduces its parameter footprint.

As presented in Table~\ref{tab:spaj15b_model_results}, xLSTM achieves lower validation perplexity relative to all other evaluated methods. For additional experimental details, see Appendix~\ref{sec:appExpComparison}. Scaling dynamics depicted in Figure~\ref{fig:temp_sclaw15B} further suggest that xLSTM retains its advantage as model size increases.

\paragraph{Ablation Studies.}
As presented in Table~\ref{tab:spaj15b_model_results} and Figure~\ref{fig:temp_sclaw15B}, xLSTM attains strong language modeling performance when trained on a dataset comprising 15B tokens from SlimPajama. To systematically evaluate the impact of individual modifications from standard LSTM to xLSTM, we incrementally transition a conventional LSTM model towards the xLSTM architecture. The process begins with incorporating LSTM layers into a residual backbone that utilizes pre-LayerNorm. Next, the architecture is augmented with a post up-projection block. The final step introduces exponential gating mechanisms together with matrix memory.

The results are shown in Table~\ref{tab:ablstudies} (top). 

Ablation experiments identify both exponential gating and matrix memory as key factors in the observed enhancement of model performance. Given the critical role of gating in RNNs and State Space Models, we further investigate alternative gating strategies. The results, shown in Table~\ref{tab:ablstudies} (bottom), indicate that allowing every gate to be learnable and responsive to input consistently yields better outcomes. Further detailed analysis regarding each backbone component can be found in Appendix~\ref{sec:appAblDetails}.

\subsection{xLSTM as Large Language Model}
\label{sec:ExpLanguage}

We conclude our investigation by conducting extensive language modeling experiments at scale, exploring xLSTM's viability as a large language model (LLM). To this end, we increase our dataset and train xLSTM using 300B tokens sourced from SlimPajama, mirroring the token count employed in Mamba~\citep{Gu:24arxiv} and Griffin~\citep{De:24arxiv}.

For comparison, we evaluate xLSTM alongside RWKV-4, Llama, and Mamba—each representing their respective methodological categories as outlined in Section~\ref{sec:ExpComparison}. RWKV-4 is chosen as the RNN exemplar, given that optimal precision settings for RWKV-5, RWKV-6, and HGRN2 were established only after commencing the 300B token experiments (refer to Appendix~\ref{sec:appGenTrainProc}).

We train a variety of model sizes (125M, 350M, 760M, 1.3B), examine length extrapolation capabilities, and measure validation set performance. Additionally, downstream task effectiveness is evaluated, language modeling proficiency is tested across 571 domains from the PALOMA benchmark, and scaling law properties are further analyzed.

\paragraph{Sequence Length Extrapolation.}
\label{sec:spaj300B_ctx_extrapolation}
We begin by assessing the sequence length extrapolation capabilities of 1.3B-parameter versions of xLSTM, RWKV-4, Llama, and Mamba. Although all models are trained with a context length of 2048, they are evaluated on significantly longer sequences, extending up to 16384 tokens. The outcomes of these evaluations are displayed in Figure~\ref{fig:spaj300B_seq_extrapolation}. Notably, xLSTM stands out by preserving lower perplexity scores on longer sequences relative to the other approaches.

\paragraph{Validation Perplexity and Downstream Tasks.}
\label{sec:spaj_downstream_eval}
Next, we analyze all scales of xLSTM, RWKV-4, Llama, and Mamba models for next token prediction on the SlimPajama validation set. Additionally, we consider their effectiveness across downstream benchmarks that probe for common sense reasoning.

The validation perplexities for various approaches are shown in the third column of Table~\ref{tab:lmeval}. Across all evaluated model sizes, xLSTM[1:0] and xLSTM[7:1] yield the lowest perplexity on the validation set, indicating superior performance. Additional columns in Table~\ref{tab:lmeval} summarize results on downstream tasks. xLSTM generally outperforms alternative methods for most tasks and sizes; however, Mamba exhibits the top performance for certain cases within the ARC task. Further information can be found in Appendix~\ref{sec:appExpLanguage}.

\paragraph{Performance on PALOMA Language Tasks.} Additionally, we evaluate the next token prediction capabilities of xLSTM, RWKV-4, Llama, and Mamba across all available model sizes using the PALOMA language tasks~\citep{Magnusson:23arxivshort}. Perplexity serves as the metric for assessing performance on next token prediction over a diverse set of 571 text domains, including sources such as nytimes.com and the Reddit community r/depression.

The results are presented in Table~\ref{tab:ppleval}, where perplexity is reported for both broad language modeling (first seven columns) and specific domain benchmarks (final five columns). On these tasks, xLSTM[1:0] consistently outperforms xLSTM[7:1].

Specifically, xLSTM[1:0] achieves lower perplexity compared to Mamba in 568 out of 571 domains (99.5\%), surpasses Llama in 486 out of 571 cases (85.1\%), and outperforms RWKV-4 in 570 out of 571 domains (99.8\%). Further analysis is available in Appendix~\ref{sec:appExpLanguage}.

\paragraph{Scaling Laws.} Next, we examine the power-law scaling phenomenon, which provides a means to predict model performance as the parameter count increases~\citep{Kaplan:20,Brown:20short}. 
The scaling characteristics are illustrated in Figure~\ref{fig:temp_sclaw300B}. Although each model demonstrates broadly comparable scaling trends, they differ in their baseline performance. RWKV-4 yields the lowest results, followed by Llama and Mamba. xLSTM outperforms Mamba by a margin analogous to Mamba's advantage over Llama. These scaling dynamics suggest that as model capacity grows, xLSTM maintains a superior trajectory relative to both Transformer architectures and State-Space models.

\textbf{Generation Times and Maximal Throughput.} In Figure~\ref{fig:inference_time_generation}, we present measurements of text generation latency, alongside an assessment of maximum throughput in Figure~\ref{fig:inference_time_decoding_throughput} (left), focusing on our xLSTM models at the 1.3B parameter scale. The evaluation includes a comparison with Mamba, Llama, and RWKV implementations of similar scale from HuggingFace, ensuring that the Llama model also benefits from a static key-value cache. Notably, attempts to compile the full Transformer model cache or utilize \texttt{torch.compile} with Mamba during the experimentation period were unsuccessful. For text generation benchmarking, each model is evaluated under batch size 1 with a pre-fill of 16, a scenario that is particularly advantageous for the Transformer architecture. As illustrated in Figure~\ref{fig:inference_time_generation}, xLSTM and other recurrent baselines (Mamba and RWKV-4) demonstrate linear complexity with respect to input size, whereas Llama exhibits the expected quadratic pattern. In terms of decoding throughput, we vary both batch size and prefill values for Llama. The analysis depicted in Figure~\ref{fig:inference_time_decoding_throughput} (right) reveals that xLSTM is able to process substantially larger batch sizes compared to Llama, attributed to its fixed memory usage, thereby achieving the highest throughput among the models considered.

\section{Limitations}

(i) Unlike mLSTM, sLSTM's approach to memory mixing restricts the use of parallelizable computation, rendering efficient parallel execution unattainable. Nonetheless, we succeeded in creating an accelerated CUDA kernel for sLSTM that is currently less than twice as slow compared to our parallel mLSTM implementation. (ii) The CUDA kernels employed in mLSTM lack optimization, resulting in the present implementation performing at about one-fourth the speed of FlashAttention or the scan method applied in Mamba. There is potential to develop faster CUDA kernels modeled after FlashAttention. (iii) The matrix memory mechanism in mLSTM presents substantial computational demands since it involves manipulating $d \times d$ matrices. However, its memory update and retrieval steps are parameter-free and can be carried out concurrently with standard matrix operations, causing only a modest increase in wall clock time attributable to the complex memory structure. (iv) Careful consideration is essential when initializing the forget gates. (v) Given that the matrix memory does not depend on sequence length, expanding the sequence length could overwhelm the memory for large context sizes. Despite this, context lengths up to 16k do not seem to pose a practical constraint, as discussed in Section~\ref{sec:spaj300B_ctx_extrapolation}. (vi) Owing to the considerable computational expense associated with large language models, we did not perform thorough optimization of the architecture or its hyperparameters, particularly for the bigger xLSTM variants. We believe that exhaustive optimization is crucial for xLSTM to fully realize its capabilities.

\section{Conclusion}
We have provided a partial response to our primary inquiry: What are the outcomes of scaling LSTM architectures to billions of parameters in the context of language modeling? Up to this point, we can state: ``We reach performance comparable to modern approaches such as Transformers and State Space Models''. By incorporating exponential gating and memory mixing along with a reimagined memory structure, we advanced LSTM into the xLSTM framework. xLSTM models consistently yield strong results in language modeling, rivalling state-of-the-art alternatives like Transformers and State Space Models. Scaling law analyses reveal that expanding xLSTM to larger sizes positions it as a formidable contender against contemporary Large Language Models based on Transformers. Furthermore, xLSTM holds promise for significantly influencing related domains such as Reinforcement Learning, Time Series Prediction, and physical system modeling.

\end{document}